% -------------------------------------------------------------------
\chapter*{Glosario de Términos}
\label{glosario}
\addcontentsline{toc}{chapter}{Glosario de Términos}

\begin{description}
% -------------------------------------------------------------------
% TCROC
% -------------------------------------------------------------------
\item[\textbf{TCROC} \(\bigl(\alpha_T\bigr)\)] 
\textbf{Tasa de Cambio Relativa Óptima Conmutada}:  
Modelo que estima la tasa de cambio relativa entre valores consecutivos de una serie temporal mediante la minimización de diferencias cuadradas. Su parámetro principal es:
\begin{itemize}
  \item \(\alpha_T\): Tasa que minimiza la suma de diferencias cuadradas entre valores adyacentes, representando la variación relativa óptima en la serie.
\end{itemize}

% -------------------------------------------------------------------
% TCROCM
% -------------------------------------------------------------------
\item[\textbf{TCROCM} \(\bigl(\alpha_T, W\bigr)\)] 
\textbf{Tasa de Cambio Relativa Óptima Conmutada Móvil}:  
Extensión de la TCROC que incorpora una ventana móvil para análisis localizado, capturando dinámicas recientes en series temporales. Sus parámetros son:
\begin{itemize}
  \item \(\alpha_T\): Tasa que optimiza la relación entre valores consecutivos dentro de la ventana móvil.
  \item \(W\): Tamaño de la ventana móvil que define el subconjunto de datos analizado.
\end{itemize}

% -------------------------------------------------------------------
% ETCROC
% -------------------------------------------------------------------
\item[\textbf{ETCROC} \(\bigl(\alpha_T, \lambda\bigr)\)] 
\textbf{Tasa de Cambio Relativa Óptima Conmutada con Decaimiento Exponencial}:  
Modelo que pondera observaciones recientes mediante un factor de decaimiento exponencial, mejorando la sensibilidad a cambios temporales. Sus parámetros son:
\begin{itemize}
  \item \(\alpha_T\): Tasa que optimiza la relación entre valores consecutivos, ajustada por ponderaciones exponenciales.
  \item \(\lambda\): Factor de decaimiento exponencial (\(\lambda \in (0,1]\)) que prioriza datos recientes.
\end{itemize}

% -------------------------------------------------------------------
% ETCROCM
% -------------------------------------------------------------------
\item[\textbf{ETCROCM} \(\bigl(\alpha_T, \lambda, W\bigr)\)] 
\textbf{Tasa de Cambio Relativa Óptima Conmutada con Decaimiento Exponencial y Ventana Móvil}:  
Combina la ventana móvil de TCROCM y el decaimiento exponencial de ETCROC para un análisis localizado y ponderado de series temporales. Sus parámetros son:
\begin{itemize}
  \item \(\alpha_T\): Tasa que optimiza la relación entre valores consecutivos en la ventana móvil, ponderada exponencialmente.
  \item \(\lambda\): Factor de decaimiento exponencial para priorizar datos recientes.
  \item \(W\): Tamaño de la ventana móvil que segmenta la serie en subconjuntos.
\end{itemize}

% -------------------------------------------------------------------
% TCROC Polinómica
% -------------------------------------------------------------------
\item[\textbf{TCROC Polinómica} \(\bigl(\alpha_T, \lambda, W, p\bigr)\)] 
\textbf{Tasa de Cambio Relativa Óptima Conmutada Polinómica}:  
Extensión de la ETCROCM que modela relaciones no lineales mediante polinomios de grado \(p\), adecuada para dinámicas complejas en series económicas y financieras. Sus parámetros son:
\begin{itemize}
  \item \(\alpha_T\): Tasa de cambio relativa promedio, ajustada por términos polinómicos.
  \item \(\lambda\): Factor de decaimiento exponencial para ponderar datos recientes.
  \item \(W\): Tamaño de la ventana móvil para análisis localizado.
  \item \(p\): Grado del polinomio que define la relación no lineal entre valores consecutivos.
\end{itemize}

% -------------------------------------------------------------------
% TCROC-Markov
% -------------------------------------------------------------------
\item[\textbf{TCROC-Markov} \(\bigl(\alpha_t, K\bigr)\)] 
\textbf{Tasa de Cambio Relativa Óptima Conmutada con Cadenas de Markov}:  
Extensión de la TCROC que modela tasas de cambio como un proceso de Markov con \(K\) estados, capturando transiciones entre regímenes en series temporales. Sus parámetros son:
\begin{itemize}
  \item \(\alpha_t\): Tasa de cambio relativa en el tiempo \(t\), modelada como variable de estado en una cadena de Markov.
  \item \(K\): Número de estados discretos en el modelo de Markov.
\end{itemize}

% -------------------------------------------------------------------
% ARIMA
% -------------------------------------------------------------------
\item[\textbf{ARIMA} \(\bigl(p, d, q\bigr)\)] 
\textbf{AutoRegressive Integrated Moving Average}:  
Modelo clásico de series temporales que combina componentes autorregresivos, de diferenciación y de promedio móvil para modelar dinámicas temporales. Sus parámetros son:
\begin{itemize}
  \item \(p\): Orden del componente autorregresivo (\emph{AutoRegressive}).
  \item \(d\): Grado de diferenciación (\emph{Integrated}) para lograr estacionariedad.
  \item \(q\): Orden del componente de promedio móvil (\emph{Moving Average}).
\end{itemize}
En su variante estacional (SARIMA), se incluyen parámetros \((P, D, Q)\) para modelar estacionalidad.

% -------------------------------------------------------------------
% Holt-Winters
% -------------------------------------------------------------------
\item[\textbf{Holt-Winters}] 
\textbf{Método de Holt-Winters}:  
Modelo que descompone series temporales en nivel, tendencia y estacionalidad, adecuado para datos con patrones estacionales pronunciados.

% -------------------------------------------------------------------
% SMA
% -------------------------------------------------------------------
\item[\textbf{SMA}] 
\textbf{Media Móvil Simple}:  
Calcula el promedio aritmético de los valores en una ventana móvil, suavizando fluctuaciones a corto plazo en la serie.

% -------------------------------------------------------------------
% EMA
% -------------------------------------------------------------------
\item[\textbf{EMA}] 
\textbf{Media Móvil Exponencial}:  
Aplica un factor de suavizado exponencial que asigna mayor peso a observaciones recientes, ideal para detectar tendencias a corto plazo.

% -------------------------------------------------------------------
% WMA
% -------------------------------------------------------------------
\item[\textbf{WMA}] 
\textbf{Media Móvil Ponderada}:  
Asigna pesos diferenciados a los valores en una ventana móvil, priorizando datos más recientes para reflejar cambios recientes.

% -------------------------------------------------------------------
% EWMA
% -------------------------------------------------------------------
\item[\textbf{EWMA}] 
\textbf{Media Móvil Exponencialmente Ponderada}:  
Similar a la EMA, utiliza ponderaciones que decrecen exponencialmente para datos más antiguos, enfatizando dinámicas recientes.

% -------------------------------------------------------------------
% Mínimos Cuadrados
% -------------------------------------------------------------------
\item[\textbf{Mínimos Cuadrados}] 
Método estadístico que estima parámetros minimizando la suma de los cuadrados de las diferencias entre valores observados y predichos.

% -------------------------------------------------------------------
% MLE
% -------------------------------------------------------------------
\item[\textbf{MLE}] 
\textbf{Estimación de Máxima Verosimilitud}:  
Técnica estadística que selecciona parámetros maximizando la probabilidad de observar los datos muestrales, usada en modelos probabilísticos.

% -------------------------------------------------------------------
% Regresión No Lineal
% -------------------------------------------------------------------
\item[\textbf{Regresión No Lineal}] 
Método para modelar relaciones no lineales entre variables dependientes e independientes, utilizando funciones no lineales ajustadas a los datos.

% -------------------------------------------------------------------
% Método Gauss-Newton
% -------------------------------------------------------------------
\item[\textbf{Método Gauss-Newton}] 
Algoritmo iterativo que resuelve problemas de mínimos cuadrados no lineales mediante linearizaciones sucesivas, optimizando parámetros en modelos no lineales.

\end{description}
% -------------------------------------------------------------------
% AIC
% -------------------------------------------------------------------
\begin{description}
\item[\textbf{AIC}] 
\textbf{Criterio de Información de Akaike}:  
Método para seleccionar modelos estadísticos, equilibrando el ajuste del modelo y su complejidad mediante la penalización del número de parámetros.

\end{description}

% -------------------------------------------------------------------
% Cadena de Markov
% -------------------------------------------------------------------
\begin{description}
\item[\textbf{Cadena de Markov}] 
\textbf{Proceso Estocástico de Markov}: 
Modelo matemático que describe una secuencia de eventos donde la probabilidad de cada evento futuro depende únicamente del estado actual. Sus componentes clave son:
\begin{itemize}
    \item Espacio de estados (\(S\)): El conjunto finito de todos los posibles regímenes.
    \item Matriz de transición (\(P\)): Matriz que contiene las probabilidades de pasar de un estado a otro.
\end{itemize}
\end{description}

% -------------------------------------------------------------------
% Codificación One-Hot
% -------------------------------------------------------------------
\begin{description}
\item[\textbf{Codificación One-Hot}] 
Proceso para convertir variables categóricas (como los estados $s_1$, $s_2$) en vectores binarios (e.g., $[1, 0, 0, 0]^T$), un formato necesario para la estimación matricial.
\end{description}

% -------------------------------------------------------------------
% Distribución Estacionaria
% -------------------------------------------------------------------
\begin{description}
\item[\textbf{Distribución Estacionaria} \(\bigl(\pi\bigr)\)]
\textbf{Distribución Estacionaria}: 
Vector de probabilidad que describe la distribución a largo plazo de una cadena de Markov, representando el equilibrio del sistema.
\end{description}

% -------------------------------------------------------------------
% Lipschitz Continuidad
% -------------------------------------------------------------------
\begin{description}
\item[\textbf{Lipschitz Continuidad}]
\textbf{Continuidad de Lipschitz}:
Propiedad matemática de una función que acota su tasa de cambio.
En esta tesis, el estimador $\alpha_t$ satisface una cota de
Lipschitz (Lema~4.4), pero la función de cuantización $\pi$ es
discontinua y no posee esta propiedad.
\end{description}

% -------------------------------------------------------------------
% SSRC
% -------------------------------------------------------------------
\begin{description}
\item[\textbf{SSRC}]
\textbf{Stochastically Structured Reservoir Computing}:
Extensión del modelo TCROC-Markov que incorpora una capa de
computación de reservorio con estructura estocástica,
permitiendo la captura de dinámicas no lineales.
\end{description}

% -------------------------------------------------------------------
% ESN
% -------------------------------------------------------------------
\begin{description}
\item[\textbf{ESN}]
\textbf{Echo State Network}:
Tipo de red neuronal recurrente donde la matriz de pesos
internos (reservorio) es fija y solo se entrena la capa
de salida, lo que simplifica el entrenamiento.
\end{description}

% -------------------------------------------------------------------
% K-Means
% -------------------------------------------------------------------
\begin{description}
\item[\textbf{K-Means}]
Algoritmo de agrupamiento que particiona $n$ observaciones en
$K$ clústeres minimizando la inercia intra-grupo.  En esta
tesis, se usa para discretizar la serie continua $\{\alpha_t\}$
en $K$ estados de régimen económico.
\end{description}

% -------------------------------------------------------------------
% Matriz de Transición
% -------------------------------------------------------------------
\begin{description}
\item[\textbf{Matriz de Transición} \(\bigl(P\bigr)\)]
\textbf{Matriz de Transición Estocástica}:
Matriz cuadrada que define las probabilidades de una cadena de Markov. Su elemento principal es:
\begin{itemize}
    \item \(P_{ij}\): Probabilidad de transitar al estado \(s_i\) desde el estado \(s_j\).
\end{itemize}
\end{description}

% -------------------------------------------------------------------
% NNLS
% -------------------------------------------------------------------
\begin{description}
\item[\textbf{NNLS}]
\textbf{Mínimos Cuadrados No Negativos}:
Algoritmo de optimización que resuelve un problema de mínimos cuadrados con la restricción de que los coeficientes deben ser no-negativos, asegurando probabilidades válidas.
\end{description}

% -------------------------------------------------------------------
% Producto de Kronecker
% -------------------------------------------------------------------
\begin{description}
\item[\textbf{Producto de Kronecker} \(\bigl(\otimes\bigr)\)]
\textbf{Producto de Kronecker}:
Operación entre matrices que resulta en una matriz de bloques más grande, utilizada para linealizar ecuaciones matriciales.
\end{description}

% -------------------------------------------------------------------
% Vectorización
% -------------------------------------------------------------------
\begin{description}
\item[\textbf{Vectorización} \(\bigl(\text{vec}(\cdot)\bigr)\)]
\textbf{Operador de Vectorización}:
Transformación lineal que convierte una matriz en un vector columna apilando sus columnas.
\end{description}

% -------------------------------------------------------------------
% phi-mixing
% -------------------------------------------------------------------
\begin{description}
\item[\textbf{\(\phi\)-mixing}]
\textbf{Mezcla Fuerte (\(\phi\)-mixing)}:
Propiedad estadística que mide la tasa a la que las observaciones lejanas en el tiempo se vuelven independientes, usada para probar propiedades asintóticas.
\end{description}
